\documentclass{beamer}

\pdfmapfile{+sansmathaccent.map}


\mode<presentation>
{
	\usetheme{Warsaw} % or try Darmstadt, Madrid, Warsaw, Rochester, CambridgeUS, ...
	\usecolortheme{seahorse} % or try seahorse, beaver, crane, wolverine, ...
	\usefonttheme{serif}  % or try serif, structurebold, ...
	\setbeamertemplate{navigation symbols}{}
	\setbeamertemplate{caption}[numbered]
} 


%%%%%%%%%%%%%%%%%%%%%%%%%%%%
% itemize settings


%%%%%%%%%%%%%%%%%%%%%%%%%%%%
% itemize settings

\definecolor{myhotpink}{RGB}{255, 80, 200}
\definecolor{mywarmpink}{RGB}{255, 60, 160}
\definecolor{mylightpink}{RGB}{255, 80, 200}
\definecolor{mypink}{RGB}{255, 30, 80}
\definecolor{mydarkpink}{RGB}{155, 25, 60}

\definecolor{mypaleblue}{RGB}{240, 240, 255}
\definecolor{mylightblue}{RGB}{120, 150, 255}
\definecolor{myblue}{RGB}{90, 90, 255}
\definecolor{mygblue}{RGB}{70, 110, 240}
\definecolor{mycourseblue}{RGB}{120, 140, 255}
\definecolor{mycourselightblue}{RGB}{245, 250, 255}
\definecolor{mydarkblue}{RGB}{0, 0, 180}
\definecolor{myblackblue}{RGB}{40, 40, 120}

\definecolor{mygreen}{RGB}{0, 200, 0}
\definecolor{mygreen2}{RGB}{245, 255, 230}
\definecolor{mydarkgreen}{RGB}{0, 120, 0}
\definecolor{mydarkgreen}{RGB}{0, 120, 0}


\definecolor{mygray}{gray}{0.8}
\definecolor{mydarkgray}{RGB}{80, 80, 160}
\definecolor{mylightgray}{RGB}{160, 160, 160}

\definecolor{mydarkred}{RGB}{160, 30, 30}
\definecolor{mylightred}{RGB}{255, 150, 150}
\definecolor{myred}{RGB}{200, 110, 110}
\definecolor{myblackred}{RGB}{120, 40, 40}

\definecolor{mypink}{RGB}{255, 30, 80}
\definecolor{myhotpink}{RGB}{255, 80, 200}
\definecolor{mywarmpink}{RGB}{255, 60, 160}
\definecolor{mylightpink}{RGB}{255, 80, 200}
\definecolor{mydarkpink}{RGB}{155, 25, 60}
\definecolor{mywhitepink}{RGB}{255, 240, 240}

\definecolor{mydarkcolor}{RGB}{60, 25, 155}
\definecolor{mylightcolor}{RGB}{130, 180, 250}



\setbeamertemplate{itemize items}[default]

\setbeamertemplate{itemize item}{\color{myblackblue}$\blacksquare$}
\setbeamertemplate{itemize subitem}{\color{mygblue}$\blacktriangleright$}
\setbeamertemplate{itemize subsubitem}{\color{mygray}$\blacksquare$}

\setbeamercolor{palette quaternary}{fg=white,bg=mycourseblue}
\setbeamercolor{titlelike}{parent=palette quaternary}

\setbeamercolor{palette quaternary2}{fg=black,bg=mycourselightblue}
\setbeamercolor{frametitle}{parent=palette quaternary2}

\setbeamerfont{frametitle}{size=\Large,series=\scshape}
\setbeamerfont{framesubtitle}{size=\normalsize,series=\upshape}





%%%%%%%%%%%%%%%%%%%%%%%%%%%%
% block settings

\setbeamercolor{block title}{bg=red!30,fg=black}

\setbeamercolor*{block title example}{bg=mygreen!40!white,fg=black}

\setbeamercolor*{block body example}{fg= black, bg= mygreen2}


%%%%%%%%%%%%%%%%%%%%%%%%%%%%
% URL settings
\hypersetup{
	colorlinks=true,
	linkcolor=blue,
	filecolor=blue,      
	urlcolor=blue,
}

%%%%%%%%%%%%%%%%%%%%%%%%%%

\renewcommand{\familydefault}{\rmdefault}


\newcommand{\mycolor}[1] {
\setbeamercolor{normal text}{fg={#1}} % Recolors standard body text
\setbeamercolor{structure}{fg={#1}}   % Recolors bullets and structure items
\usebeamercolor[fg]{normal text}      % Applies the normal text color immediately
}


\usepackage{amsmath}
\usepackage{mathtools}

\usepackage{cancel}
\usepackage{centernot}

\usepackage{changepage}
\usepackage{subcaption}

\usepackage{qrcode}

\newcommand{\bo}[1] {\mathbf{#1}}
\newcommand{\R}{\mathbb{R}} 
\newcommand{\T}{^\top}     



\newcommand{\mydate}{Fall 2026}

\newcommand{\mygit}{\textcolor{blue}{\href{https://github.com/SergeiSa/Advanced-Control-2026}{github.com/SergeiSa/Advanced-Control-2026}}}

\newcommand{\myqr}{ \textcolor{black}{\qrcode[height=1.5in]{https://github.com/SergeiSa/Advanced-Control-2026}}
}

\newcommand{\myqrframe}{
	\begin{frame}
		\centerline{Lecture slides are available via Github:}
		\bigskip
		\centerline{\mygit}
		\bigskip
		\myqr
	\end{frame}
}


\newcommand{\bref}[2] {\textcolor{blue}{\href{#1}{#2}}}

%%%%%%%%%%%%%%%%%%%%%%%%%%%%
% code settings

\usepackage{listings}
\usepackage{color}
% \definecolor{mygreen}{rgb}{0,0.6,0}
% \definecolor{mygray}{rgb}{0.5,0.5,0.5}
\definecolor{mymauve}{rgb}{0.58,0,0.82}
\lstset{ 
	backgroundcolor=\color{white},   % choose the background color; you must add \usepackage{color} or \usepackage{xcolor}; should come as last argument
	basicstyle=\footnotesize,        % the size of the fonts that are used for the code
	breakatwhitespace=false,         % sets if automatic breaks should only happen at whitespace
	breaklines=true,                 % sets automatic line breaking
	captionpos=b,                    % sets the caption-position to bottom
	commentstyle=\color{mygreen},    % comment style
	deletekeywords={...},            % if you want to delete keywords from the given language
	escapeinside={\%*}{*)},          % if you want to add LaTeX within your code
	extendedchars=true,              % lets you use non-ASCII characters; for 8-bits encodings only, does not work with UTF-8
	firstnumber=0000,                % start line enumeration with line 0000
	frame=single,	                   % adds a frame around the code
	keepspaces=true,                 % keeps spaces in text, useful for keeping indentation of code (possibly needs columns=flexible)
	keywordstyle=\color{blue},       % keyword style
	language=Octave,                 % the language of the code
	morekeywords={*,...},            % if you want to add more keywords to the set
	numbers=left,                    % where to put the line-numbers; possible values are (none, left, right)
	numbersep=5pt,                   % how far the line-numbers are from the code
	numberstyle=\tiny\color{mygray}, % the style that is used for the line-numbers
	rulecolor=\color{black},         % if not set, the frame-color may be changed on line-breaks within not-black text (e.g. comments (green here))
	showspaces=false,                % show spaces everywhere adding particular underscores; it overrides 'showstringspaces'
	showstringspaces=false,          % underline spaces within strings only
	showtabs=false,                  % show tabs within strings adding particular underscores
	stepnumber=2,                    % the step between two line-numbers. If it's 1, each line will be numbered
	stringstyle=\color{mymauve},     % string literal style
	tabsize=2,	                   % sets default tabsize to 2 spaces
	title=\lstname                   % show the filename of files included with \lstinputlisting; also try caption instead of title
}


%%%%%%%%%%%%%%%%%%%%%%%%%%%%
% URL settings
\hypersetup{
	colorlinks=false,
	linkcolor=blue,
	filecolor=blue,      
	urlcolor=blue,
}

%%%%%%%%%%%%%%%%%%%%%%%%%%

%%%%%%%%%%%%%%%%%%%%%%%%%%%%
% tikz settings

\usepackage{tikz}
\tikzset{every picture/.style={line width=0.75pt}}
\usetikzlibrary{shapes.geometric, decorations.markings, positioning, patterns, calc, arrows.meta}


\title{Adjoint Method, Sensitivity Propagation}
\subtitle{Advanced Control}
\author{by Sergei Savin}
\centering
\date{\mydate}



\begin{document}
\maketitle


\begin{frame}{Content}
\textbf{Adjoint Method}
\begin{itemize}
	\item Lagrangian and Hamiltonian derivation
	\item Costate equations and final gradient
	\item Comparison with Pontryagin's Principle
	\item Connection to reverse-mode AD
\end{itemize}

\textbf{Sensitivity Propagation Method}
\begin{itemize}
	\item Sensitivity ODE derivation
	\item Cost gradient computation
	\item Connection to forward-mode AD
\end{itemize}
\end{frame}



%------------------------------------------------
\begin{frame}{Direct Shooting: Problem Statement}
	\begin{flushleft}
		Direct shooting trajectory planning can be represented as
		the following optimization problem:
		
		\[
		\boxed{
			\begin{aligned}
				\min_{c}\quad
				J(c)&=\Phi(x(t_f))+
				\int_0^{t_f}\ell(x,u,t)\,dt\\
				\text{s.t.}\quad
				\dot{x}&=f(x,u,t),\\
				u&=u(c,t),\\
				x(0)&=x_0.
			\end{aligned}
		}
		\]
		
		\bigskip
		
		Here, $c$ parametrizes the control input, while
		$x(t)$ is determined by integrating the dynamics. $\Phi$ is the terminal cost and $\ell$ is the running cost.
	\end{flushleft}
\end{frame}

%------------------------------------------------
\begin{frame}{Gradients}
	\begin{flushleft}
		To find parameters $c$ that minimize $J(c)$,
		we need to compute the gradient $\frac{dJ}{dc}$.
		
		\bigskip
		
		Computing it explicitly is difficult because:
		\begin{itemize}
			\item The state $x(t)$ depends implicitly on $c$
			through the system dynamics.
			\item Computing $dx(t)/dc$ requires differentiating
			the ODE solution with respect to every parameter.
			\item The number of sensitivity equations grows
			with the number of parameters in $c$.
		\end{itemize}
		
		\bigskip
		
		We will derive a method that avoids computing
		$dx(t)/dc$ explicitly.
	\end{flushleft}
\end{frame}

%------------------------------------------------
\begin{frame}{Gradients Through a Lagrangian}
	\begin{flushleft}
		We introduce a Lagrangian to incorporate the equality
		constraint describing the system dynamics:
		
		\begin{align}
			\boxed{
				\begin{aligned}
					\mathcal{L}(c,\lambda)
					={}&\Phi(x(t_f))\\
					&+\int_0^{t_f}
					\left[
					\ell(x,u,t)+
					\lambda^\top(f(x,u,t)-\dot{x})
					\right]dt.
				\end{aligned}
			}
		\end{align}
		
		Here, $\lambda(t)$ is a time-dependent Lagrange multiplier,
		also called the \textbf{costate}.
		
		\bigskip
		
		On the set $\dot{x}=f(x,u,t)$, which includes every
		admissible trajectory,
		\[
		\boxed{\mathcal{L}(c,\lambda)=J(c).}
		\]
	\end{flushleft}
\end{frame}

%------------------------------------------------
\begin{frame}{Simplification of the Lagrangian}
	\begin{flushleft}
		A key step is to eliminate $\dot{x}$ using integration
		by parts.
		
		\begin{theorem}[Integration by parts]
			\[
			\int_a^b u^\top\dot{v}\,dt
			=
			\left[u^\top v\right]_a^b
			-\int_a^b \dot{u}^\top v\,dt.
			\]
		\end{theorem}
		
		Applying this formula gives
		\[
		-\int_0^{t_f}\lambda^\top\dot{x}\,dt
		=
		-\lambda(t_f)^\top x(t_f)
		+\lambda(0)^\top x_0
		+\int_0^{t_f}\dot{\lambda}^\top x\,dt.
		\]
		
		Thus, the Lagrangian becomes
		\[
		\boxed{
			\begin{aligned}
				\mathcal{L}
				={}&\Phi(x(t_f))-\lambda(t_f)^\top x(t_f)+\lambda(0)^\top x_0\\
				&+\int_0^{t_f}
				\left[\ell+\lambda^\top f+
				\dot{\lambda}^\top x\right]dt.
			\end{aligned}
		}
		\]
	\end{flushleft}
\end{frame}

%------------------------------------------------
\begin{frame}{Gradient with Respect to Parameters}
	\begin{flushleft}
		Differentiate $\mathcal{L}$ with respect to $c$,
		treating $\lambda(t)$ as fixed:
		
		\[
		\begin{aligned}
			\frac{dJ}{dc}
			={}&
			\left(\frac{d\Phi}{dx}
			-\lambda(t_f)^\top\right)
			\frac{dx(t_f)}{dc}\\
			&+\lambda(0)^\top\frac{dx_0}{dc}\\
			&+\int_0^{t_f}
			\left(
			\frac{d\ell}{dx}
			+\lambda^\top\frac{df}{dx}
			+\dot{\lambda}^\top
			\right)\frac{dx}{dc}\,dt\\
			&+\int_0^{t_f}
			\left(
			\frac{d\ell}{du}
			+\lambda^\top\frac{df}{du}
			\right)\frac{du}{dc}\,dt.
		\end{aligned}
		\]
		
		\bigskip
		
		Since $x(0)=x_0$ is independent of $c$,
		\[
		\frac{dx_0}{dc}=0.
		\]
	\end{flushleft}
\end{frame}

%------------------------------------------------
\begin{frame}{Simplification via the Hamiltonian}
	\begin{flushleft}
		Define the Hamiltonian:
		\[
		\boxed{
			H(x,u,\lambda,t)=\ell(x,u,t)+\lambda^\top f(x,u,t).
		}
		\]
		
		Its derivatives are
		\[
		\begin{aligned}
			\frac{dH}{dx}
			&=\frac{d\ell}{dx}+\lambda^\top\frac{df}{dx},\\
			\frac{dH}{du}
			&=\frac{d\ell}{du}+\lambda^\top\frac{df}{du}.
		\end{aligned}
		\]
		
		\bigskip
		
		The gradient can now be written as
		\[
		\boxed{
			\begin{aligned}
				\frac{dJ}{dc}
				={}&
				\left(\frac{d\Phi}{dx}-\lambda(t_f)^\top\right)
				\frac{dx(t_f)}{dc}\\
				&+\int_0^{t_f}
				\left(\frac{dH}{dx}+\dot{\lambda}^\top\right)
				\frac{dx}{dc}\,dt\\
				&+\int_0^{t_f}\frac{dH}{du}\frac{du}{dc}\,dt.
			\end{aligned}
		}
		\]
	\end{flushleft}
\end{frame}

%------------------------------------------------
\begin{frame}{Eliminating State Gradients}
	\begin{flushleft}
		We choose $\lambda(t)$ to cancel all terms containing
		$dx/dc$.
		
		\bigskip
		
		Canceling the terminal term gives
		\[
		\boxed{
			\lambda(t_f)=\left(\frac{d\Phi}{dx}\right)^\top.
		}
		\]
		
		\bigskip
		
		Canceling the integral term gives
		\[
		\boxed{
			\dot{\lambda}
			=-\left(\frac{dH}{dx}\right)^\top.
		}
		\]
		
		\bigskip
		
		This is an ODE solved \textbf{backward in time},
		using $\lambda(t_f)$ as its terminal condition.
	\end{flushleft}
\end{frame}

%------------------------------------------------
\begin{frame}{Gradients: Final Form}
	\begin{flushleft}
		With the costate chosen as above, all terms involving
		$dx/dc$ vanish.
		
		\bigskip
		
		We obtain the final expression for the gradient:
		\[
		\boxed{
			\frac{dJ}{dc}
			=
			\int_0^{t_f}
			\frac{dH}{du}\frac{du}{dc}\,dt.
		}
		\]
		
		\bigskip
		
		Equivalently,
		\[
		\boxed{
			\frac{dJ}{dc}
			=
			\int_0^{t_f}
			\left(
			\frac{d\ell}{du}
			+\lambda^\top\frac{df}{du}
			\right)
			\frac{du}{dc}\,dt.
		}
		\]
		
		The state sensitivities $dx/dc$ are no longer needed.
	\end{flushleft}
\end{frame}

%------------------------------------------------
\begin{frame}{Description of the Adjoint Method}
	\begin{flushleft}
		The adjoint method computes the gradient in four steps:
		
		\begin{enumerate}
			\item \textbf{Forward pass:}
			Integrate $\dot{x}=f(x,u,t)$ and store the trajectory.
			
			\item \textbf{Backward pass:}
			Integrate $\dot{\lambda}=-(dH/dx)^\top$ backward,
			starting from $\lambda(t_f)=(d\Phi/dx)^\top$.
			
			\item \textbf{Gradient evaluation:}
			Compute $\displaystyle
			\frac{dJ}{dc}=\int_0^{t_f}
			\frac{dH}{du}\frac{du}{dc}\,dt$.
			
			\item \textbf{Solver update:}
			Update $c$ using a gradient-based optimizer
			and repeat until convergence.
		\end{enumerate}
		
		\bigskip
		
		The cost of the backward pass does not grow with
		the number of control parameters.
	\end{flushleft}
\end{frame}

%------------------------------------------------
\begin{frame}{Comparison with Pontryagin's Principle}
	\begin{flushleft}
		
		\textbf{Pontryagin's minimum principle} gives necessary
		conditions for an optimal control trajectory. For an optimal solution, it introduces a Hamiltonian
		and costate satisfying
		\[
		\dot{\lambda}=-\left(\frac{dH}{dx}\right)^\top,
		\qquad
		\lambda(t_f)=\left(\frac{d\Phi}{dx}\right)^\top.
		\]
		
		For an unconstrained optimal control, it also requires
		\[
		\frac{dH}{du}=0
		\]
		along an optimal trajectory.
		
		\textbf{The key difference:}
		\begin{itemize}
			\item The adjoint method computes the gradient
			$dJ/dc$ for a chosen parameterized control.
			\item Pontryagin's principle provides necessary
			optimality conditions for the optimal control itself.
		\end{itemize}
		
	\end{flushleft}
\end{frame}
%------------------------------------------------



%------------------------------------------------
\begin{frame}{Reverse-Mode Automatic Differentiation vs.\ the Adjoint Method}
	\begin{flushleft}
		
		\textbf{Discrete case (shooting transcription):} Reverse-mode AD on the unrolled integration loop
		\emph{is} the discrete adjoint method.
		
		\bigskip
		
		\textbf{Limitations of the AD implementation:} Reverse-Mode AD builds a computational graph (often large) that could require substantial memory to store intermediate values.
		
	\end{flushleft}
\end{frame}
%------------------------------------------------




%------------------------------------------------
\begin{frame}{Direct Shooting: Sensitivity}
	\begin{flushleft}
		The state trajectory depends on the control parameters
		through the dynamics and the input $u(c,t)$.
		
		For each parameter $c_j$, define the state sensitivity
		\[
		S_j(t) := \frac{dx(t)}{dc_j}.
		\]
		
		\bigskip
		
		\textbf{Idea:} Derive an ODE governing the evolution of
		$S_j(t)$. Once the sensitivities are known, we can compute
		the gradient $\frac{dJ}{dc}$ using the chain rule.
	\end{flushleft}
\end{frame}

%------------------------------------------------
\begin{frame}{Direct Shooting: Differentiating the Dynamics}
	\begin{flushleft}
		Start with the dynamics
		\[
		\dot{x} = f(x,u,t),
		\qquad u=u(c,t).
		\]
		
		Differentiate both sides with respect to $c_j$:
		\[
		\frac{d}{dc_j}\dot{x}
		=
		\frac{df}{dx}\frac{dx}{dc_j}
		+
		\frac{df}{du}\frac{du}{dc_j}.
		\]
		
		Since $c_j$ is independent of time, differentiation
		with respect to $c_j$ commutes with differentiation
		with respect to $t$:
		\[
		\frac{d}{dc_j}\dot{x}
		=
		\frac{d}{dt}\frac{dx}{dc_j}.
		\]
		
		Defining $S_j := \frac{dx}{dc_j},
		\quad
		U_j := \frac{du}{dc_j}$, 
		we obtain:
		\[
		\boxed{
			\dot{S}_j
			=
			\frac{df}{dx}S_j
			+
			\frac{df}{du}U_j
		}
		\]
	\end{flushleft}
\end{frame}

%------------------------------------------------
\begin{frame}{Direct Shooting: Sensitivity ODE}
	\begin{flushleft}
		Collect the sensitivities with respect to all $N$
		control parameters into matrices:
		\[
		S =
		\begin{bmatrix}
			S_1 & \cdots & S_N
		\end{bmatrix}
		=
		\frac{dx}{dc},
		\qquad
		U_c = \frac{du}{dc}.
		\]
		
		Concatenating the individual sensitivity equations gives
		\[
		\boxed{
			\dot{S}(t)
			=
			\frac{df}{dx}S(t)
			+
			\frac{df}{du}U_c(t)
		}
		\]
		
		\bigskip
		
		Because $x(0)=x_0$ is independent of $c$,
		\[
		\boxed{S(0)=0.}
		\]
		
		\textbf{Numerical procedure:} Integrate the state and
		sensitivity ODEs forward together. Evaluate $\frac{df}{dx}$
		and $\frac{df}{du}$ along the state and input trajectories.
	\end{flushleft}
\end{frame}

%------------------------------------------------
\begin{frame}{Direct Shooting: Cost Gradient}
	\begin{flushleft}
		The objective is
		\[
		J(c)=\Phi(x(t_f))
		+\int_0^{t_f}\ell(x,u,t)\,dt.
		\]
		
		Apply the chain rule, using
		$S(t)=\frac{dx}{dc}$ and $U_c(t)=\frac{du}{dc}$:
		\[
		\boxed{
			\begin{aligned}
				\frac{dJ}{dc}
				={}&
				\frac{d\Phi}{dx}S(t_f)\\
				&+
				\int_0^{t_f}
				\left(
				\frac{d\ell}{dx}S(t)
				+
				\frac{d\ell}{du}U_c(t)
				\right)dt.
			\end{aligned}
		}
		\]
		
		\bigskip
		
		Thus, the state and sensitivity ODEs provide the
		trajectory and its derivatives, from which we compute
		$\frac{dJ}{dc}$.
	\end{flushleft}
\end{frame}



%------------------------------------------------
\begin{frame}{Sensitivity Propagation vs Forward-Mode AD}
	\begin{flushleft}
		
		Discrete-time sensitivity propagation is the \emph{forward-mode automatic differentiation}:

		\bigskip

		For a discrete system $x_{k+1}=F_k(x_k,c)$, both Sensitivity Propagation and Forward-Mode AD are represented as:
		
		\[ \boxed{ S_{k+1} = \frac{dF_k}{dx_k}S_k + \frac{dF_k}{dc}, \qquad S_0=0. } \]
		
		
	\end{flushleft}
\end{frame}





\begin{frame}{Literature}
	
	\begin{itemize}
		
		
		\item Russ Tedrake -- MIT:
		\bref{https://underactuated.mit.edu/trajopt.html}
		{Underactuated Robotics: Trajectory Optimization, ch. 10.}
		
	\end{itemize}
	
\end{frame}



\myqrframe




\end{document}
